A análise e classificação de dados complexos são desafios centrais em Inteligência Computacional, especialmente diante de conjuntos de dados reais que apresentam alta dimensionalidade, múltiplas classes e estruturas não triviais. Métodos de agrupamento (clustering) desempenham papel fundamental na descoberta de padrões, segmentação e análise exploratória em diversas áreas, como ciência de dados, engenharia, saúde e ciências sociais.

Neste trabalho, diferentes algoritmos de agrupamento, incluindo KMeans, Fuzzy C-Means, Gustafson-Kessel e Agglomerativo, são implementados, documentados e comparados utilizando os conjuntos de dados Adult e Dry Bean. Cada método é explorado em notebooks didáticos, com etapas bem definidas: importação de bibliotecas, carregamento e pré-processamento dos dados, implementação dos modelos, aplicação, visualização dos resultados e análise crítica.

A comparação entre os algoritmos é realizada a partir de métricas como acurácia, precisão, revocação, especificidade, F1-score e análise da matriz de confusão, permitindo avaliar o desempenho de cada abordagem frente aos desafios impostos pelos dados. Para facilitar a interpretação dos agrupamentos e dos centros dos clusters, utiliza-se a técnica de Análise de Componentes Principais (PCA) para projeção dos dados em duas dimensões, possibilitando uma visualização clara da separação entre os grupos encontrados.

A abordagem adotada busca não apenas comparar o desempenho dos algoritmos, mas também fornecer uma base didática sólida para compreensão dos conceitos envolvidos, destacando as vantagens e limitações de cada método no contexto dos dados analisados.